% ================================================================
%  conteudo/resultados.tex
% ================================================================

\section{Resultados e Discussão}

Após o período de incubação de seis dias a 30~°C, observou-se o desenvolvimento de colônias distintas nas placas inoculadas tanto com o material da superfície externa das uvas quanto com a polpa macerada. A Figura~1 apresenta o conjunto das seis placas de Petri ao término do experimento, permitindo uma visão geral do perfil de crescimento microbiano obtido nas diferentes condições de amostragem.

\begin{figure}[h]
\centering
\caption{Placas de Petri após seis dias de incubação a 30~°C}
\includegraphics[width=0.75\textwidth]{image/todas_placas.png}
\fonte{Os autores (2026)}
\end{figure}

Dentre as colônias desenvolvidas, identificaram-se morfologias compatíveis com leveduras, caracterizadas por colônias de coloração creme a esbranquiçada, superfície lisa, bordas regulares e consistência pastosa, conforme descrito por Vieira e Fernandes (2012). A Figura~2 evidencia, em detalhe, uma placa na qual se observa o crescimento característico de colônias de leveduras isoladas pelo método de esgotamento por estrias.

\begin{figure}[h]
\centering
\caption{Detalhe de placa com colônias de leveduras isoladas}
\includegraphics[width=0.65\textwidth]{image/placa_com_levedura.png}
\fonte{Os autores (2026)}
\end{figure}

O crescimento observado está em conformidade com o esperado para a microbiota natural de uvas, cuja superfície é conhecida por abrigar populações diversificadas de leveduras nativas (BERES, 2013). A predominância de morfologias leveduriformes nas placas pode ser atribuída à seletividade do meio de cultura empregado, que, associada ao pH ácido característico das uvas, favorece o desenvolvimento fúngico em detrimento de bactérias (EEEP, 2012). A técnica de esgotamento por estrias mostrou-se eficiente para a separação e individualização das colônias, corroborando as observações de Câmara (2023) acerca da adequação desse método para o isolamento de microrganismos a partir de amostras mistas de origem vegetal.
